\chapter{Challenges, Limitations and Future Work}
\label{chapter:discussion}

\section{Current Limitations}
\label{sec:limitations}

Unfortunately this QLBM has a crucial limitation: each time step requires measurements and a new initialization of the state.
It is especially oppressive as the encoding (which was excluded from the above complexity) is linear in the number of lattice sites.

It is still unclear if the higher error for complex velocity fields is due to the algorithm itself or our implementation. Further investigation is required.

The algorithm requires quantum state tomography at each iteration which is a significant limitation for practical usage.
There is active research in this area, such as the work of Kocherla et al. \cite{kocherla2024fully}, in which they can defer measurement until the end of the simulation.
Combining their approach with our implementation is an interesting avenue.

Both the initialization procedure and the collision operator are highly inefficient, relying on Qiskit's transpiler to generate complex and deep circuits.

\todo{REVISION: reword future directions}
\section{Future Directions}
\label{sec:future}

There are more recent developments that try to solve these problems.
For instance, Kocherla et al. proposed a QLBM that can delay the need for measurements until the end of the simulation \cite{kocherla2024fully}.


The algorithm could be extended to simulate Navier-Stokes flows by using the stream-vorticity formulation, as proposed by Budinski \cite{ljubomir2022quantum}.

There are plenty of directions for future work regarding boundary conditions:
\begin{itemize}
    \item Investigating the kinds of boundary conditions that can be expressed using this method, and whether or not all post-streaming conditions can be rewritten as pre-streaming ones.
    \item Extending the method to support constants in the boundary equations, allowing for source and sink terms.
    \item Extending the method to support non-local boundary conditions that depend on multiple nodes.
    \item Testing the method in higher-dimensional simulations.
    \item Exploring other methods for the LCU decomposition that would avoid SVD.
\end{itemize}
